\documentclass[../../main.tex]{subfiles}
\graphicspath{{../../image/}{../../figure/}}
\begin{document}

\section{稳定性估计}\label{section:ch15-稳定性估计}

本节给出稳定性估计\eqref{align:ch15-1515}、\eqref{align:ch15-1516}与\eqref{align:ch15-1517}证明的梗概。

由于缩放 $U(x,t)=U_\mu(\mu x,\mu t)$，\eqref{align:ch15-1513}、\eqref{align:ch15-1512}解的估计\eqref{align:ch15-1516}与\eqref{align:ch15-1517}分别等价于\eqref{align:ch15-1521}、\eqref{align:ch15-1522}解的估计
\begin{align}
\|U(\cdot,t)-U(\cdot,\tau)\|_{L^1(-\infty,\infty)}\leqslant b\bigl(|t-\tau|+|\sqrt{t}-\sqrt{\tau}|\bigr),\qquad\tau,t\in[0,\infty),
\label{align:ch15-1581}
\end{align}
\begin{align}
\|U(\cdot,t)-\bar U(\cdot,t)\|_{L^1(-\infty,\infty)}\leqslant a\|U_{0\mu}(\cdot)-\bar U_{0\mu}(\cdot)\|_{L^1(-\infty,\infty)},\qquad t\in(0,\infty).
\label{align:ch15-1582}
\end{align}

估计\eqref{align:ch15-1581}通过在 $(\tau,t)$ 上积分不等式
\begin{align}
\|\partial_tU(\cdot,t)\|_{L^1(-\infty,\infty)}\leqslant b\left(1+\frac1{2\sqrt{t}}\right),\qquad t\in(0,\infty),
\label{align:ch15-1583}
\end{align}
而得到；该不等式由\eqref{align:ch15-1521}、并借助\eqref{align:ch15-1523}、\eqref{align:ch15-1531}与\eqref{align:ch15-1524}推出。

估计\eqref{align:ch15-1582}由下述同伦论证建立。我们有
\begin{align}
U(x,t)-\bar U(x,t)=\int_0^1\frac{\mathrm{d}}{\mathrm{d}\xi}U_\xi(x,t)\,\mathrm{d}\xi,
\label{align:ch15-1584}
\end{align}
其中 $U_\xi$ 表示\eqref{align:ch15-1521}以 $\xi\bar U_{0\mu}+(1-\xi)U_{0\mu}$ 为初值的解。“切”向量
\begin{align}
W_\xi(x,t)=\frac{\mathrm{d}}{\mathrm{d}\xi}U_\xi(x,t)
\label{align:ch15-1585}
\end{align}
是线性化方程
\begin{align}
\partial_tW_\xi(x,t)+\partial_x\left[A(U_\xi(x,t))W_\xi(x,t)\right]=\partial_x^2W_\xi(x,t),
\label{align:ch15-1586}
\end{align}
在初值
\begin{align}
W_\xi(\cdot,0)=\bar U_{0\mu}(\cdot)-U_{0\mu}(\cdot)
\label{align:ch15-1587}
\end{align}
下的解。

方程\eqref{align:ch15-1586}与\eqref{align:ch15-1521}解的导数 $\partial_xU$ 所满足的方程非常相似，因而可用前几节所用方法处理。\cref{section:ch15-注记}所引文献中的分析表明：当 $\|\partial_xU_\xi(\cdot,t)\|_{L^2(-\infty,\infty)}<\delta_0$ 在 $(0,\infty)$ 上成立时，存在常数 $a>1$，使得对任意 $\delta>0$，$\|W_\xi(\cdot,0)\|_{L^1(-\infty,\infty)}<\delta/a$ 蕴含 $\|W_\xi(\cdot,t)\|_{L^1(-\infty,\infty)}<\delta$，对一切 $t\in(0,\infty)$。由于\eqref{align:ch15-1586}是线性的，上述断言等价于
\begin{align}
\|W_\xi(\cdot,t)\|_{L^1(-\infty,\infty)}\leqslant a\|W_\xi(\cdot,0)\|_{L^1(-\infty,\infty)},\qquad0<t<\infty.
\label{align:ch15-1588}
\end{align}

联合\eqref{align:ch15-1588}、\eqref{align:ch15-1584}、\eqref{align:ch15-1585}与\eqref{align:ch15-1587}，即得\eqref{align:ch15-1582}，从而建立\eqref{align:ch15-1517}。

其余估计\eqref{align:ch15-1515}是\eqref{align:ch15-1517}的直接推论。事实上，对两个解 $U_\mu(x,t)$ 与 $\bar U_\mu(x,t)=U_\mu(x+h,t)$ 应用\eqref{align:ch15-1517}，其相应初值为 $U_0(x)$ 与 $\bar U_0(x)=U_0(x+h)$，把所得不等式乘以 $h^{-1}$ 再令 $h\to0$，即得\eqref{align:ch15-1515}。

由消失黏性法构造的\eqref{align:ch15-1511}解具有有限传播速度性质。事实上，利用 Green 函数的性质可证：当 $U_0$ 与 $\bar U_0$ 在区间 $(y,z)$ 内重合时（此时对 $x\in(y/\mu,z/\mu)$ 有 $U_{0\mu}(x)=\bar U_{0\mu}(x)$），相应的\eqref{align:ch15-1521}、\eqref{align:ch15-1522}解 $U$ 与 $\bar U$ 满足
\begin{align}
|U(x,t)-\bar U(x,t)|\leqslant c\|U_0(\cdot)-\bar U_0(\cdot)\|_{L^\infty(-\infty,\infty)}\left\{\exp\left(vt-x+\frac y\mu\right)+\exp\left(vt+x-\frac z\mu\right)\right\}
\label{align:ch15-1589}
\end{align}
其中 $c$、$v$ 为正常数，对一切 $(x,t)\in(-\infty,\infty)\times(0,\infty)$ 成立。作重标度 $(x,t)\mapsto(x/\mu,t/\mu)$ 回到 $U_\mu$、$\bar U_\mu$，可得：初值为 $U_0$、$\bar U_0$ 的两个解 $U=\lim_{k\to\infty}U_{\mu_k}$ 与 $\bar U=\lim_{k\to\infty}\bar U_{\mu_k}$，对一切 $x\in(y+vt,z-vt)$ 必有重合。

由此可知，代替\eqref{align:ch15-15110}与\eqref{align:ch15-1518}，我们有更精细的估计
\begin{align}
\int_y^z|U(x,t)-\bar U(x)|\,dx\leqslant a\int_{y-vt}^{z+vt}|U_0(x)-\bar U_0(x)|\,dx,
\label{align:ch15-15810}
\end{align}
\begin{align}
TV_{(y,z)}U(\cdot,t)\leqslant a\,TV_{(y-vt,z+vt)}U_0(\cdot),
\label{align:ch15-15811}
\end{align}
对任意 $-\infty\leqslant y<z\leqslant\infty$ 成立。

进而证明：有限传播速度性质与稳定性估计\eqref{align:ch15-1518}联合，蕴含由消失黏性法构造的\eqref{align:ch15-1511}、\eqref{align:ch15-1512}任一解 $U$ 满足 Tame Oscillation Condition 14.10.1。

由于\eqref{align:ch15-1511}的解在空间与时间平移下保持，只需在原点 $x=0$，$t=0$ 处验证 (14.10.3)。
\begin{align}
\bar U_0(x)=
\begin{cases}
U_0(-\lambda h+)&-\infty<x\leqslant-\lambda h,\\
U_0(x)&-\lambda h<x<\lambda h,\\
U_0(\lambda h-)&\lambda h<x<\infty
\end{cases}
\label{align:ch15-15812}
\end{align}
为初值的解 $\bar U$。此时有 $TV_{(-\infty,\infty)}\bar U_0(\cdot)=TV_{(-\lambda h,\lambda h)}U_0(\cdot)$，$\bar U(0\pm,h)=U(0\pm,h)$，$\bar U(\infty,h)=U_0(\lambda h-)$。因此，由\eqref{align:ch15-1518}，
\begin{align}
|U(0\pm,h)-U_0(0\pm)|&\leqslant|\bar U(0\pm,h)-\bar U(\infty,h)|+|U_0(\lambda h-)-U_0(0\pm)|\nonumber\\
&\leqslant(a+1)TV_{(-\lambda h,\lambda h)}U_0(\cdot),
\label{align:ch15-15813}
\end{align}
这就建立了 (14.10.3)。

\end{document}
